\chapter{Estação de Monitoramento IoT}
\label{cap:projeto-final}

Este capítulo final reúne, em um único projeto, praticamente tudo o que foi apresentado nesta apostila: tipos de dados, estruturas de controle, funções, vetores, \code{structs} e \code{enums}, além de GPIO, ADC, comunicação serial, interrupções, Watchdog Timer, Wi-Fi e atualizações OTA. O objetivo é consolidar esses conceitos através de um firmware completo e funcional — não apenas um exemplo isolado de cada periférico.

\section{Especificação do projeto}

Uma \textbf{estação de monitoramento} de luminosidade ambiente: o \ESP{} lê continuamente um sensor LDR, mantém uma média móvel das últimas leituras para suavizar ruído, registra o estado do ambiente pela porta serial e acende um LED de alerta sempre que o ambiente estiver escuro. Um botão alterna entre modo \textbf{normal} (só registra alertas) e modo \textbf{verboso} (registra cada leitura) — sem nunca bloquear o laço principal com \code{delay}. Além disso, a estação se conecta ao Wi-Fi para reportar alertas a um servidor, aceita atualizações de firmware pela rede (OTA), e é protegida por um Watchdog Timer contra travamentos.

\begin{nota}
O \emph{deep sleep} do \cref{cap:energia} não entra no firmware a seguir: ele é incompatível com o botão de interrupção e com a resposta imediata que este projeto busca (durante o \emph{deep sleep}, o processador está desligado e não pode reagir a nada além das fontes de despertar configuradas, nem manter a conexão Wi-Fi). Já o \emph{light sleep} — a opção mais viável para um projeto como este, por preservar o estado do programa e continuar acordando pelo próprio botão — é discutido na seção final deste capítulo.
\end{nota}

\subsection{Materiais}

\begin{itemize}
    \item 1 placa de desenvolvimento \ESP{};
    \item 1 LED (com resistor limitador em série, \textasciitilde220~$\Omega$), como indicador de alerta;
    \item 1 sensor LDR, em um divisor de tensão com um resistor fixo (\textasciitilde10~k$\Omega$);
    \item 1 botão \emph{push-button};
    \item Protoboard e jumpers.
\end{itemize}

\subsection{Ligações sugeridas}

\begin{itemize}
    \item LED de alerta $\rightarrow$ GPIO2 (através do resistor limitador, ao GND);
    \item Ponto médio do divisor de tensão do LDR $\rightarrow$ GPIO34 (ADC1);
    \item Botão $\rightarrow$ GPIO15 e GND (usando o \emph{pull-up} interno, como no \cref{cap:gpio}).
\end{itemize}

\section{Projetando a estrutura do estado}

Antes de escrever o laço principal, vale modelar o \textbf{estado} do sistema com os tipos vistos no \cref{cap:structs} — uma prática de organização que se torna cada vez mais valiosa à medida que um firmware cresce:

\begin{codigoc}{Modelando o estado do sistema}
enum ModoOperacao {
    MODO_NORMAL,
    MODO_VERBOSO,
};

struct EstadoEstacao {
    ModoOperacao modo;
    int ultima_leitura;
    float media_movel;
    bool alerta_ativo;
};
\end{codigoc}

Agrupar essas variáveis em uma única \code{struct}, em vez de espalhá-las como variáveis globais soltas, deixa explícito o que compõe ``o estado do sistema'' — e facilita, por exemplo, passar esse estado inteiro para uma função de log, caso o projeto cresça no futuro.

\section{Firmware completo}

\begin{codigoc}{src/main.cpp — estação de monitoramento IoT}
// Bibliotecas do framework Arduino/ESP32 usadas neste projeto:
#include <WiFi.h>          // conexao a rede Wi-Fi (cap. 21)
#include <HTTPClient.h>    // requisicoes HTTP para reportar alertas (cap. 21)
#include <ArduinoOTA.h>    // atualizacao de firmware pela rede (cap. 22)
#include <esp_task_wdt.h>  // Watchdog Timer, protege contra travamentos (cap. 19)

// Pinos usados no projeto (ver secao "Ligacoes sugeridas")
#define PINO_LED    2    // LED de alerta
#define PINO_LDR    34   // sensor de luminosidade (entrada ADC1)
#define PINO_BOTAO  15   // botao que alterna o modo de operacao

// Parametros de comportamento do firmware, centralizados aqui para
// facilitar ajustes sem precisar caçar "numeros magicos" pelo codigo
#define TAMANHO_HISTORICO   10      // quantas leituras entram na media movel
#define LIMIAR_ESCURO       25.0f   // abaixo disso, ambiente "escuro" (%)
#define INTERVALO_LEITURA   500UL   // ms entre leituras do sensor
#define WDT_TIMEOUT         5       // segundos ate o reset do watchdog

// Credenciais e destino da rede -- em um projeto real, evite deixar
// senhas fixas no codigo-fonte (considere um arquivo separado e ignorado
// pelo controle de versao, como discutido no cap. 23 do apendice de Git)
#define WIFI_SSID     "MinhaRede"
#define WIFI_PASSWORD "minha_senha"
#define URL_SERVIDOR  "http://exemplo.com/api/alertas"

// Os dois modos de operacao alternados pelo botao (cap. 10)
enum ModoOperacao {
    MODO_NORMAL,    // registra apenas quando um alerta comeca
    MODO_VERBOSO,   // registra toda leitura do sensor
};

// Todo o estado "vivo" do sistema, reunido em uma unica struct (cap. 10)
// em vez de espalhado em variaveis globais soltas
struct EstadoEstacao {
    ModoOperacao modo;     // modo atual (NORMAL ou VERBOSO)
    int ultima_leitura;    // valor bruto (0-4095) da ultima leitura do ADC
    float media_movel;     // media das ultimas TAMANHO_HISTORICO leituras
    bool alerta_ativo;     // true enquanto o ambiente esta "escuro"
};

// Estado inicial: modo normal, nenhuma leitura ainda, sem alerta
EstadoEstacao estado = {MODO_NORMAL, 0, 0.0f, false};

// Buffer circular (cap. 8) com as ultimas leituras, usado para calcular
// a media movel que suaviza o ruido natural do sensor LDR
int historico[TAMANHO_HISTORICO];
int indice_historico = 0;   // proxima posicao do buffer a ser sobrescrita

// 'volatile' avisa o compilador que esta variavel pode mudar a qualquer
// momento por causa de uma interrupcao -- sem isso, o compilador poderia
// otimizar o codigo assumindo (erradamente) que ela nunca muda em loop()
volatile bool evento_botao = false;
unsigned long ultima_leitura_ms = 0;   // controla a temporizacao sem delay()

// Rotina de interrupcao (ISR) do botao: deve ser curta e rapida, apenas
// sinalizando o evento -- todo o trabalho de fato acontece em loop().
// IRAM_ATTR garante que o codigo da ISR fique carregado na RAM interna,
// para nao depender de uma leitura da flash externa durante a interrupcao
void IRAM_ATTR botao_isr() {
    evento_botao = true;
}

// Calcula a media aritmetica simples do buffer circular 'historico'
float calcular_media(void) {
    long soma = 0;   // 'long' evita overflow ao somar ate 10 leituras de 12 bits
    for (int i = 0; i < TAMANHO_HISTORICO; i++) {
        soma += historico[i];
    }
    // cast para float antes de dividir: divisao inteira arredondaria a media
    return (float) soma / TAMANHO_HISTORICO;
}

// Conecta ao Wi-Fi, bloqueando ate a conexao ser estabelecida -- aceitavel
// aqui porque so acontece uma vez, dentro de setup(), antes do laco comecar
void conectar_wifi(void) {
    WiFi.begin(WIFI_SSID, WIFI_PASSWORD);
    Serial.print("Conectando ao Wi-Fi");

    while (WiFi.status() != WL_CONNECTED) {
        delay(500);
        Serial.print(".");   // feedback visual de que o firmware nao travou
    }

    Serial.println();
    Serial.print("Conectado! IP: ");
    Serial.println(WiFi.localIP());
}

// Envia o percentual de luminosidade atual para o servidor via HTTP POST,
// no formato JSON, sempre que um novo alerta comeca
void enviar_alerta(float luminosidade_pct) {
    // sem Wi-Fi conectado, nao ha para onde enviar -- desiste sem travar
    if (WiFi.status() != WL_CONNECTED) {
        return;
    }

    HTTPClient http;
    http.begin(URL_SERVIDOR);
    http.addHeader("Content-Type", "application/json");

    // snprintf monta o corpo da requisicao em um buffer de tamanho fixo,
    // evitando a fragmentacao de memoria associada a concatenar String
    // (cap. 18) -- sizeof(corpo) impede que a escrita ultrapasse o buffer
    char corpo[64];
    snprintf(corpo, sizeof(corpo), "{\"luminosidade\": %.1f}", luminosidade_pct);

    int codigo = http.POST(corpo);
    Serial.print("[HTTP] Alerta enviado, resposta: ");
    Serial.println(codigo);   // codigo de status HTTP (200, 404, etc.)

    http.end();   // libera os recursos da requisicao
}

void setup() {
    Serial.begin(115200);

    pinMode(PINO_LED, OUTPUT);
    pinMode(PINO_BOTAO, INPUT_PULLUP);
    // dispara botao_isr() na borda de descida (botao pressionado, pull-up)
    attachInterrupt(digitalPinToInterrupt(PINO_BOTAO), botao_isr, FALLING);

    // zera o historico para que a primeira media movel nao misture
    // leituras validas com "lixo" de memoria nao inicializada
    for (int i = 0; i < TAMANHO_HISTORICO; i++) {
        historico[i] = 0;
    }

    conectar_wifi();   // bloqueia aqui ate a rede estar pronta

    // deixa a estacao pronta para receber uma atualizacao de firmware
    // pela rede a qualquer momento, sem precisar de cabo (cap. 22)
    ArduinoOTA.setHostname("estacao-iot");
    ArduinoOTA.begin();

    // segundo argumento 'true': trava (panic) o sistema em vez de so
    // resetar, o que ajuda a diagnosticar travamentos durante o desenvolvimento
    esp_task_wdt_init(WDT_TIMEOUT, true);
    esp_task_wdt_add(NULL);   // registra a tarefa atual (loop) no watchdog

    Serial.println("Estacao de monitoramento iniciada em modo NORMAL.");
}

void loop() {
    ArduinoOTA.handle();   // fica pronta para uma atualizacao a qualquer momento

    // 1) verifica se o botao foi pressionado (sem bloquear o laco)
    if (evento_botao) {
        evento_botao = false;   // consome o evento, para nao tratar de novo
        // alterna entre os dois modos com o operador ternario (cap. 5)
        estado.modo = (estado.modo == MODO_NORMAL) ? MODO_VERBOSO : MODO_NORMAL;

        Serial.print("Modo alterado para: ");
        Serial.println(estado.modo == MODO_NORMAL ? "NORMAL" : "VERBOSO");
    }

    // 2) le o sensor periodicamente, sem usar delay() -- temporizacao
    //    por millis(), a mesma tecnica usada nos exemplos do cap. 19
    unsigned long agora = millis();
    if (agora - ultima_leitura_ms >= INTERVALO_LEITURA) {
        ultima_leitura_ms = agora;   // reinicia a contagem do intervalo

        // leitura bruta do ADC (0 a 4095, resolucao de 12 bits no ESP32)
        estado.ultima_leitura = analogRead(PINO_LDR);

        // grava a leitura na posicao atual do buffer circular e avanca
        // o indice, voltando a 0 automaticamente com o operador '%'
        historico[indice_historico] = estado.ultima_leitura;
        indice_historico = (indice_historico + 1) % TAMANHO_HISTORICO;

        estado.media_movel = calcular_media();
        // converte o valor bruto do ADC (0-4095) para um percentual (0-100)
        float luminosidade_pct = (estado.media_movel / 4095.0f) * 100.0f;

        // guarda o estado anterior do alerta para detectar a TRANSICAO
        // de "sem alerta" para "em alerta" logo abaixo
        bool estava_em_alerta = estado.alerta_ativo;
        estado.alerta_ativo = (luminosidade_pct < LIMIAR_ESCURO);

        digitalWrite(PINO_LED, estado.alerta_ativo);   // acende/apaga o LED

        // no modo verboso, registra toda leitura; no modo normal, so
        // os alertas sao registrados (ver bloco de alerta logo abaixo)
        if (estado.modo == MODO_VERBOSO) {
            Serial.print("[LEITURA] bruta=");
            Serial.print(estado.ultima_leitura);
            Serial.print(" media=");
            Serial.print(estado.media_movel);
            Serial.print(" luminosidade=");
            Serial.print(luminosidade_pct);
            Serial.println("%");
        }

        // envia o alerta apenas no instante em que ele comeca (nao a cada
        // leitura) -- e por isso que 'estava_em_alerta' foi guardado acima
        if (estado.alerta_ativo && !estava_em_alerta) {
            Serial.println("[ALERTA] Ambiente escuro detectado!");
            enviar_alerta(luminosidade_pct);
        }
    }

    esp_task_wdt_reset();   // "alimenta" o watchdog a cada volta saudavel do laco
}
\end{codigoc}

\section{Como o projeto se encaixa}

\begin{itemize}
    \item A \code{struct EstadoEstacao} e a \code{enum ModoOperacao} (\cref{cap:structs}) organizam o estado do sistema de forma clara, em vez de espalhá-lo em variáveis globais soltas;
    \item O vetor \code{historico} (\cref{cap:arranjos}) implementa um \textbf{buffer circular} para a média móvel, suavizando o ruído natural do sensor, exatamente como praticado no \cref{cap:adc};
    \item A leitura do LDR usa a \textbf{ADC} (\cref{cap:adc}) para medir luminosidade de forma contínua, convertida de valor bruto para porcentagem com um \emph{cast} (\cref{cap:tipos});
    \item O botão é lido por \textbf{interrupção} (\cref{cap:interrupcoes}), sinalizando o laço principal por meio de uma variável \code{volatile}, sem nunca bloquear a leitura contínua do sensor;
    \item A leitura periódica do sensor usa a técnica de \textbf{temporização por \code{millis()}} (\cref{cap:interrupcoes}), mantendo \code{loop()} sempre livre para reagir ao botão e às atualizações OTA a qualquer momento;
    \item Um \textbf{Watchdog Timer} (\cref{sec:watchdog}) protege o laço principal como um todo, reiniciando o dispositivo automaticamente caso algo trave a execução: a cada volta saudável de \code{loop()}, o firmware chama \code{esp\_task\_wdt\_reset()} para ``alimentá-lo'';
    \item A estação se conecta à rede via \textbf{Wi-Fi} (\cref{cap:wifi}) e reporta cada novo alerta a um servidor por \textbf{HTTP}, usando \code{snprintf} para montar o corpo da requisição sem a fragmentação de memória associada a concatenar \code{String};
    \item O firmware aceita \textbf{atualizações OTA} (\cref{cap:ota}) através de \code{ArduinoOTA.handle()}, chamado a cada volta do laço — a estação pode receber correções e melhorias sem nunca precisar ser desconectada e levada até um computador;
    \item Todo o comportamento é registrado pela \textbf{porta serial} (\cref{cap:serial}) — mensagens de leitura apenas em modo verboso, e alertas sempre, independentemente do modo.
\end{itemize}

\section{Sobre economizar energia neste projeto}

Vale retomar a ressalva feita no início do capítulo: o \emph{deep sleep} do \cref{cap:energia} \textbf{não} foi incorporado a este firmware, e essa foi uma decisão deliberada, não um esquecimento. Este projeto foi especificado para ficar \emph{sempre alerta} — reagir ao botão instantaneamente, aceitar uma atualização OTA a qualquer momento, manter a conexão Wi-Fi ativa — e o \emph{deep sleep} desliga exatamente os componentes (processador, rádio) que tornam isso possível, reiniciando o firmware do zero a cada ciclo.

O \textbf{\emph{light sleep}} (\cref{sec:light-sleep}) é uma história diferente, e de fato a opção mais viável se este projeto precisasse rodar por mais tempo em uma bateria: como preserva toda a RAM e pode acordar por qualquer pino GPIO — incluindo o próprio \code{PINO\_BOTAO} já usado aqui —, ele não exige abandonar nem a estrutura do estado (\code{struct EstadoEstacao}), nem a responsividade ao botão. A forma mais direta de incorporá-lo seria substituir o \code{delay(100)} tipicamente usado em trabalhos curtos e síncronos por uma pausa de \emph{light sleep} equivalente, ou intercalar breves períodos de sono entre uma leitura do sensor e a seguinte, quando não há nenhum evento pendente para tratar.

A complicação — e o motivo de o exemplo completo desta apostila não incluir isso — é que \code{esp\_light\_sleep\_start()} interage com a conexão \textbf{Wi-Fi} e a disponibilidade do \textbf{OTA} de um jeito que exige coordenação cuidadosa com o rádio (\cref{sec:light-sleep}), algo que foge do escopo introdutório deste material. Ainda assim, entre as duas opções de baixo consumo apresentadas no \cref{cap:energia}, é o \emph{light sleep} — não o \emph{deep sleep} — que representa o próximo passo natural para tornar esta estação mais eficiente em energia sem sacrificar sua responsividade.

Se, por outro lado, o objetivo fosse um sensor remoto alimentado por bateria \emph{sem} necessidade de resposta imediata a um botão físico nem de manter o Wi-Fi permanentemente ativo, o projeto certo seria mais radical: acordar a cada poucos minutos, ler o sensor, conectar ao Wi-Fi \emph{apenas} para enviar o alerta (se houver um), e voltar ao \emph{deep sleep} — o padrão ``acordar, medir, dormir'' do \cref{cap:energia}, em vez da arquitetura de laço único sempre ativo usada aqui. Essa tensão entre \textbf{responsividade} e \textbf{autonomia de energia} é uma decisão de projeto real, que todo engenheiro de sistemas embarcados precisa fazer conscientemente — não existe uma resposta certa universal, apenas a mais adequada a cada aplicação.

\section{Para continuar praticando}

Esta apostila termina aqui, mas o aprendizado de sistemas embarcados é, acima de tudo, prático. Como próximos passos a partir deste projeto, considere:

\begin{itemize}
    \item Adicionar um segundo sensor (por exemplo, um sensor de temperatura) e expandir a \code{struct EstadoEstacao} para incluí-lo;
    \item Salvar o \code{modo} escolhido na memória flash (usando a biblioteca \code{Preferences} do Arduino para ESP32), para que a estação lembre a última configuração mesmo após ser desligada;
    \item Introduzir pausas de \emph{light sleep} (\cref{sec:light-sleep}) nos trechos realmente ociosos do laço, para reduzir o consumo sem abrir mão da resposta ao botão — o passo intermediário mais realista para esta estação;
    \item Ir além, adaptando o projeto para o padrão ``acorda, mede, dorme'' com \emph{deep sleep} do \cref{cap:energia}, caso a resposta imediata ao botão físico deixe de ser um requisito;
    \item Trocar as requisições HTTP avulsas por uma conexão MQTT persistente (\cref{cap:wifi}), caso o projeto cresça para múltiplos dispositivos reportando dados continuamente;
    \item Reescrever o projeto dividindo-o em múltiplos arquivos \code{.h}/\code{.cpp} (\cref{cap:modularizacao}) — por exemplo, separando a lógica do sensor, da rede e do botão em módulos independentes.
\end{itemize}

Cada um desses passos reutiliza e aprofunda os mesmos fundamentos de C apresentados nesta apostila — a linguagem muda pouco; o que cresce é a sua familiaridade com o hardware e o ecossistema ao redor dela.
